\scanpage{68}
\deleted{Un épinglage de $M$ s'identifie donc à une}

On peut considérer, indépendamment de la question d'épingler $J$, \deleted{celle de trouver une section \unsure{?|de|s}} \inserted{l'épinglage ``relatif'' \unsure{consistant|consistera}} en \deleted{une} la donnée d'une section du torseur $P_M$ sous $\underline{\mathrm{End}}(M_0)$, ce qui équivaut \inserted{\unsure{du|donc}} aussi à la donnée d'un morphisme
\[
\underline J\longrightarrow\deleted{\unsure{V_{\underline J}(}}M\deleted{\unsure{)|\}}} \tag*{(47)}
\]
qui relève le morphisme d'inclusion
\[
\underline J\hookrightarrow M_0
\]
et qui soit une section de $V_{\underline J}(M)$ i.e. qui soit à trace nulle (rel. à $S$).

\unsure{On a|Ceci est|Ceci} en fait automatiquement un morphisme
\[
\underline J\longrightarrow M\setminus M_0
\]
donc en composant avec \deleted{l'inclusion}
\[
(\underline J\dashrightarrow)\,M\hookrightarrow E\longrightarrow\Sigma_{\underline J}\simeq E/(1,\sigma_E)
\]
on trouve un morphisme
\[
\underline J\longrightarrow\Delta_t \tag*{(48)}
\]
et la condition de relèvement pour (47) signifie aussi que le composé de (48)
% archive sheet number: 67
